\documentclass[12pt, a4paper]{report}

% ─── Packages ────────────────────────────────────────────────────────────────
\usepackage[utf8]{inputenc}
\usepackage[T1]{fontenc}
\usepackage[english]{babel}
\usepackage{graphicx}
\usepackage{hyperref}
\usepackage{geometry}
\usepackage{setspace}
\usepackage{fancyhdr}
\usepackage{indentfirst}
\usepackage{caption}
\usepackage{subcaption}
\usepackage{amsmath}
\usepackage{amssymb}
\usepackage{url}
\usepackage{cite}
\usepackage{listings}
\usepackage{xcolor}
\usepackage{booktabs}
\usepackage{float}
\usepackage{lipsum} % Remove when filling in real content

% ─── Page geometry ───────────────────────────────────────────────────────────
\geometry{
    a4paper,
    top=2.5cm,
    bottom=2.5cm,
    left=3cm,
    right=2.5cm
}

% ─── Line spacing ────────────────────────────────────────────────────────────
\onehalfspacing

% ─── Header / footer ─────────────────────────────────────────────────────────
\setlength{\headheight}{15pt}
\pagestyle{fancy}
\fancyhf{}
\fancyhead[L]{\leftmark}
\fancyhead[R]{\thepage}
\renewcommand{\headrulewidth}{0.4pt}

% ─── Hyperlink colours ───────────────────────────────────────────────────────
\hypersetup{
    colorlinks=true,
    linkcolor=black,
    citecolor=black,
    urlcolor=blue
}

% ─── Code listing style ──────────────────────────────────────────────────────
\lstset{
    basicstyle=\ttfamily\small,
    keywordstyle=\color{blue},
    commentstyle=\color{gray},
    stringstyle=\color{red},
    breaklines=true,
    frame=single,
    numbers=left,
    numberstyle=\tiny\color{gray}
}

% ─────────────────────────────────────────────────────────────────────────────
%  DOCUMENT
% ─────────────────────────────────────────────────────────────────────────────
\begin{document}

% ─── Title page ──────────────────────────────────────────────────────────────
\begin{titlepage}
    \centering
    \vspace*{1cm}

    {\large \textbf{University Name}}\\[0.3cm]
    {\large Department of Computer Science}\\[2cm]

    % \includegraphics[width=0.3\textwidth]{logo}\\[2cm] % uncomment and replace 'logo' with your logo file

    {\LARGE \textbf{Project Title}}\\[0.5cm]
    {\large Technical Report}\\[2cm]

    {\large \textbf{Author Name}}\\[0.2cm]
    {\large Student Number: ioh123456}\\[2cm]

    {\large Supervisor: Prof. Supervisor Name}\\[0.5cm]
    {\large Co-supervisor: Prof. Co-supervisor Name}\\[2cm]

    {\large \today}

    \vfill
\end{titlepage}

% ─── Abstract ────────────────────────────────────────────────────────────────
\chapter*{Abstract}
\addcontentsline{toc}{chapter}{Abstract}

% TODO: Write the abstract (150–250 words recommended).
% Summarise: motivation, problem, approach, main results and conclusions.
\lipsum[1-2]

\vspace{1em}
\noindent\textbf{Keywords:} keyword1, keyword2, keyword3, keyword4, keyword5.

% ─── Table of contents ───────────────────────────────────────────────────────
\tableofcontents
\listoffigures     % remove if you have no figures
\listoftables      % remove if you have no tables

% ─────────────────────────────────────────────────────────────────────────────
%  CHAPTER 1 – Introduction
% ─────────────────────────────────────────────────────────────────────────────
\chapter{Introduction}
\label{chap:introduction}

\section{Context and Motivation}
% TODO: Describe the context of the project and why it is relevant.
\lipsum[3]

\section{Problem Statement}
% TODO: Clearly define the problem being addressed.
\lipsum[4]


\section{Goals and Objectives}
% TODO: List the main goals of the project.
The main goals of this project are:
\begin{itemize}
    \item Goal 1;
    \item Goal 2;
    \item Goal 3.
\end{itemize}

\section{Document Structure}
% TODO: Briefly describe what each chapter covers.
This report is organised as follows.
Chapter~\ref{chap:related_work} presents related work, tools and technologies.
Chapter~\ref{chap:requirements} describes the requirements elicitation and system architecture.
Chapter~\ref{chap:expected_results} outlines the expected results.

% ─────────────────────────────────────────────────────────────────────────────
%  CHAPTER 2 – Related Work, Tools and Technologies
% ─────────────────────────────────────────────────────────────────────────────
\chapter{Related Work, Tools and Technologies}
\label{chap:related_work}

\section{Related Work}
% TODO: Review the literature. Compare existing solutions with your approach.
\lipsum[5-6]

\subsection{Similar System / Approach A}
% TODO: Describe a related system or approach.
\lipsum[7]

\subsection{Similar System / Approach B}
% TODO: Describe another related system or approach.
\lipsum[8]

\subsection{Summary and Comparison}
% TODO: Summarise and compare the reviewed work.
\lipsum[9]

\section{Tools and Technologies}
% TODO: Describe the main tools, frameworks and technologies you plan to use.

\subsection{Technology / Framework A}
% TODO: Describe this technology and justify your choice.
\lipsum[10]

\subsection{Technology / Framework B}
% TODO: Describe this technology and justify your choice.
\lipsum[11]

\subsection{Technology / Framework C}
% TODO: Describe this technology and justify your choice.
\lipsum[12]

% ─────────────────────────────────────────────────────────────────────────────
%  CHAPTER 3 – Requirements Elicitation and Architecture
% ─────────────────────────────────────────────────────────────────────────────
\chapter{Requirements Elicitation and Architecture}
\label{chap:requirements}

\section{Requirements Elicitation}
% TODO: Explain how requirements were gathered (interviews, surveys, etc.).
Tou  a fazer isto, não mexam nisto pls :)  by pedro

\subsection{Functional Requirements}
% TODO: List and describe the functional requirements.
\begin{itemize}
    \item \textbf{FR1 – Centralização e Exposição de Serviços:} O sistema deve atuar como uma API Gateway 
                                                    centralizada para expor serviços da Universidade de Aveiro para 
                                                    ser usada nas mais diversas aplicações.
    \item \textbf{FR2 – Gestão de Identidade:} O sistema tem de implementar mecanismos de autenticação e autorização. 
                                                Nem todas as APIs podem ser usadas por todos os utilizadores. 
                                                O sistema deve garantir que apenas os utilizadores autorizados podem aceder a certas APIs.
                                                Além disso o sistemas assegurar um controlo de acessos robusto aos vários endpoints.
    \item \textbf{FR3 – Controlo e Monitorização:} A plataforma deve permitir a monitorização detalhada de todos os acessos e controlo dos mesmos.
    \item \textbf{FR4 – Disponibilização de Documentação:} O sistema deve disponibilizar documentação para utilização dos serviços na plataforma API.
    \item \textbf{FR5 – Comunicação com Sistemas Antigos:} A plataforma tem de ser capaz de se integrar com os sistemas legados da universidade.
\end{itemize}

\subsection{Non-Functional Requirements}
% TODO: List and describe the non-functional requirements (performance, security, etc.).
\begin{itemize}
    \item \textbf{NFR1} – The system shall \ldots
    \item \textbf{NFR2} – The system shall \ldots
\end{itemize}

\section{System Architecture}
% TODO: Describe the overall architecture. Use a diagram if possible.
\lipsum[14]

% Example figure placeholder – replace with your own diagram
% \begin{figure}[H]
%     \centering
%     \includegraphics[width=0.8\textwidth]{architecture_diagram}
%     \caption{System Architecture Overview}
%     \label{fig:architecture}
% \end{figure}

\subsection{Component / Module A}
% TODO: Describe this component and its responsibilities.
\lipsum[15]

\subsection{Component / Module B}
% TODO: Describe this component and its responsibilities.
\lipsum[16]

\subsection{Data Model}
% TODO: Describe the main entities and their relationships.
\lipsum[17]

% ─────────────────────────────────────────────────────────────────────────────
%  CHAPTER 4 – Expected Results
% ─────────────────────────────────────────────────────────────────────────────
\chapter{Expected Results}
\label{chap:expected_results}

% TODO: Describe what you expect to achieve at the end of the project.
\lipsum[18]

\section{Deliverables}
% TODO: List the concrete deliverables (prototype, documentation, etc.).
\begin{itemize}
    \item Deliverable 1 – Description;
    \item Deliverable 2 – Description;
    \item Deliverable 3 – Description.
\end{itemize}

\section{Evaluation Criteria}
% TODO: Define how you will evaluate the success of the project.
\lipsum[19]

\section{Project Schedule}
% TODO: Present a timeline (e.g. a Gantt chart or a simple table).
\begin{table}[H]
    \centering
    \caption{Planned project schedule}
    \label{tab:schedule}
    \begin{tabular}{clcc}
        \toprule
        \textbf{Phase} & \textbf{Task} & \textbf{Start} & \textbf{End} \\
        \midrule
        1 & Requirements elicitation      & Mar 2026 & Mar 2026 \\
        2 & Architecture design           & Mar 2026 & Apr 2026 \\
        3 & Implementation                & Apr 2026 & Jun 2026 \\
        4 & Testing and validation        & Jun 2026 & Jul 2026 \\
        5 & Report writing and submission & Jul 2026 & Jul 2026 \\
        \bottomrule
    \end{tabular}
\end{table}

% ─────────────────────────────────────────────────────────────────────────────
%  Bibliography
% ─────────────────────────────────────────────────────────────────────────────
\bibliographystyle{plain}
\bibliography{references} % create a references.bib file with your citations

% ─────────────────────────────────────────────────────────────────────────────
%  Appendices (optional for first part)
% ─────────────────────────────────────────────────────────────────────────────
% \appendix
% \chapter{Appendix A – Additional Material}

\end{document}
